\textbf{Proof.}
Write \(n=n_D=N_0K_D\) and index the donor-post patients by \(i=(m,j)\).  By \(\mathrm{lem:normalization\mbox{-}collapse}\), every donor-post event survival curve is \(s_1\).  Hence all donor-post failure times have the same hazard, denoted by \(\lambda_1\), although their censoring survivals \(G_{1,m}\) need not agree.  For patient \(i=(m,j)\), define
\[
 N_i(u)=\mathbf 1\{X_{1,m,j}\le u,\ \Delta_{1,m,j}=1\},
 \qquad
 Y_i(u)=\mathbf 1\{X_{1,m,j}\ge u\},
\]
and
\[
 M_i(u)=N_i(u)-\int_0^uY_i(v)\lambda_1(v)\,dv.       \tag{1}
\]
Control consistency and independent censoring give the compensator in (1).  Patient independence makes the processes \(M_i\) independent, mean-zero, square-integrable martingales.  Put
\[
 Y=\sum_iY_i,\qquad N=\sum_iN_i,\qquad M=\sum_iM_i,
 \qquad J_Y(u)=\mathbf 1\{Y(u)>0\}.
\]
The event hazard is continuous and therefore has no finite atoms.  Thus independent censoring and the equal donor cell sizes imply
\[
 E Y(u)
 =K_Ds_1(u)\sum_{m\in\mathcal D}G_{1,m}(u)
 =nH_D(u),
 \qquad H_D(u)\ge h_0:=s_-g_->0.                       \tag{2}
\]
All constants below depend only on the fixed constants in the model assumptions.

We first record two uniform risk-set bounds.  Since \(Y(\bar\tau)\) is a sum of independent Bernoulli variables with mean at least \(nh_0\), the exponential Markov argument for a lower Bernoulli tail gives
\[
 P\{Y(\bar\tau)<nh_0/2\}\le \exp(-c n).               \tag{3}
\]
For completeness, a uniform second-moment bound needs no identical-distribution assumption.  If \(W_i(u)=\mathbf 1\{X_i\ge u\}\), introduce independent copies and independent Rademacher signs.  Jensen's inequality, symmetrization by the copies, and the triangle inequality reduce the squared maximal centered sum to at most four times the conditional squared maximum of
\(\sum_i\epsilon_iW_i(u)\).  Conditional on the observations, as \(u\) varies these sums are reverse partial sums of a fixed ordering of the independent signs.  The elementary partial-sum maximal inequality consequently gives
\[
 E\sup_{u\in\mathcal I}
 \left|\frac{Y(u)}n-H_D(u)\right|^2\le \frac Cn.       \tag{4}
\]
Ties only select a subset of those partial sums and therefore do not change the bound.

Let \(\widehat S_D^{\mathrm{PL}}\) be the ordinary product-limit estimator formed after pooling all donor-post records, and set
\(R=\widehat S_D^{\mathrm{PL}}/s_1\).  Product integration and integration by parts yield the exact identity
\[
 R(t)-1
 =-\int_0^tR(u-)\frac{J_Y(u)}{Y(u)}\,dM(u)
  +\int_0^tR(u-)\{1-J_Y(u)\}\lambda_1(u)\,du,          \tag{5}
\]
where \(J_Y/Y\) is defined to be zero when \(Y=0\).  Indeed, substituting
\(dN=dM+Y\lambda_1du\) in the product-limit differential cancels the event-hazard drift whenever \(Y>0\); the second term in (5) is exactly the drift left after risk-set exhaustion.

Let
\[
 \sigma_n=\inf\{u:Y(u)<nh_0/2\}\wedge\bar\tau.
\]
Before \(\sigma_n\) the drift in (5) is zero.  Since
\(0\le\widehat S_D^{\mathrm{PL}}\le1\), event overlap gives
\(0\le R\le s_-^{-1}\).  The predictable isometry followed by the martingale maximal inequality therefore gives
\[
 E\sup_{t\le\sigma_n}|R(t)-1|^2
 \le 4E\int_0^{\sigma_n}
 R(u-)^2Y(u)^{-2}\,d\langle M\rangle(u)
 \le \frac Cn,                                         \tag{6}
\]
because \(d\langle M\rangle=Y\lambda_1du\),
\(Y\ge nh_0/2\) before \(\sigma_n\), and
\(\sup_u\lambda_1(u)\le L\).

Define the independent mean-zero influence processes
\[
 \phi_i^{\mathrm{pool}}(t)
 =-s_1(t)\int_0^tH_D(u)^{-1}\,dM_i(u).                 \tag{7}
\]
Subtracting \(n^{-1}\sum_i\phi_i^{\mathrm{pool}}\) from \(s_1(R-1)\) in (5) gives a martingale remainder with predictable integrand
\[
 D_n(u)=R(u-)\frac{J_Y(u)}{Y(u)}-\frac1{nH_D(u)}.       \tag{8}
\]
On the event
\(A_n=\{Y(\bar\tau)\ge nh_0/2\}\), monotonicity of \(Y\) and (2) imply
\[
 |D_n(u)|
 \le \frac{2}{nh_0}|R(u-)-1|
 +\frac{2}{nh_0^2}
   \left|\frac{Y(u)}n-H_D(u)\right|.                  \tag{9}
\]
Equations (4), (6), and (9), followed again by the predictable isometry and the martingale maximal inequality, show that
\[
 E\left[\mathbf 1_{A_n}
 \sup_{t\in\mathcal I}
 \left|s_1(t)\int_0^tD_n(u)\,dM(u)\right|^2\right]
 \le Cn^{-2}.                                          \tag{10}
\]
The same polynomial order holds on \(A_n^c\).  To see this without any unmentioned tail condition, note that successive failure risk sizes are distinct decreasing positive integers.  Consequently
\[
 \int_0^{\bar\tau}\frac{J_Y(u)}{Y(u)}\,dN(u)
 \le\sum_{k=1}^n k^{-1}\le1+\log n.                   \tag{11}
\]
The compensator parts in (5) and (8) are bounded by fixed constants, while the oracle jump sum with integrand \((nH_D)^{-1}\) is also bounded by a fixed constant.  Thus the squared supremum in (10) is at most
\(C(1+\log n)^2\) on every data set.  Combining this deterministic bound with (3) gives
\[
 E\left[\mathbf 1_{A_n^c}
 \sup_{t\in\mathcal I}
 \left|s_1(t)\int_0^tD_n(u)\,dM(u)\right|^2\right]
 \le C(1+\log n)^2e^{-cn}\le Cn^{-2}.                 \tag{12}
\]

The risk-exhaustion drift
\[
 E_n(t)=s_1(t)\int_0^tR(u-)\{1-J_Y(u)\}\lambda_1(u)\,du
\]
is uniformly bounded and vanishes unless \(Y(\bar\tau)=0\).  Equation (3), with the still smaller empty-risk event, therefore implies
\[
 \|E E_n\|_\infty\le Ce^{-cn},
 \qquad
 E\|E_n-E E_n\|_\infty^2\le Ce^{-cn}.                 \tag{13}
\]
Set
\[
 b_{D,n}^{\mathrm{PL}}=E E_n,
 \qquad
 q_{D,n}^{\mathrm{PL}}(t)
 =-s_1(t)\int_0^tD_n(u)\,dM(u)+E_n(t)-E E_n(t).        \tag{14}
\]
The stochastic integral in (14) has expectation zero because its integrand is predictable.  Hence (10)--(14) prove the exact decomposition
\[
 \widehat S_D^{\mathrm{PL}}-s_1
 =n^{-1}\sum_i\phi_i^{\mathrm{pool}}
  +b_{D,n}^{\mathrm{PL}}+q_{D,n}^{\mathrm{PL}},         \tag{15}
\]
with
\[
 \|b_{D,n}^{\mathrm{PL}}\|_\infty\le Ce^{-cn},
 \qquad E q_{D,n}^{\mathrm{PL}}=0,
 \qquad E\|q_{D,n}^{\mathrm{PL}}\|_\infty^2\le Cn^{-2}. \tag{16}
\]
Piecewise-linear interpolation is a deterministic supremum-norm contraction and fixes every grid knot.  Applying it to the three terms on the right of (15), and naming the interpolated last two terms \(b_{D,n}\) and \(q_{D,n}\), proves the expansion asserted in the statement on \(\mathcal G_{N_0,K_T,K_D}\).  The equality \(\widehat\theta_1=\widehat S_D^{\mathrm{PL}}\) on that branch and \(\theta=s_1\) follow respectively from \(\mathrm{def:fpcr\mbox{-}multiplier\mbox{-}band}\) and \(\mathrm{lem:normalization\mbox{-}collapse}\).

We next prove the two second-moment bounds.  From (2), (7), the predictable isometry, and the martingale maximal inequality,
\[
 E\left\|n^{-1}\sum_i\phi_i^{\mathrm{pool}}\right\|_\infty^2
 \le \frac4n\int_0^{\bar\tau}
       \frac{\lambda_1(u)}{H_D(u)}\,du
 \le Cn^{-1}.                                          \tag{17}
\]
Interpolation can only decrease this supremum norm.

For the empirical multiplier, condition on the records and omit the outer product-limit factor, which is bounded by one.  In failure-time order its remaining part is the partial-sum process
\[
 V^\xi(t)=\sum_i\xi_i\int_0^t\frac{J_Y(u)}{Y(u)}\,dN_i(u).
\]
The conditional partial-sum maximal inequality gives
\[
 E_\xi\sup_t|V^\xi(t)|^2
 \le4\int_0^{\bar\tau}\frac{J_Y(u)}{Y(u)^2}\,dN(u).   \tag{18}
\]
On \(A_n\), the right side is at most \(C/n\).  On \(A_n^c\), successive failure risk sizes are distinct, so it is at most
\(4\sum_{k\ge1}k^{-2}\), a fixed constant.  Averaging (18), using (3), and again using interpolation as a contraction yields
\[
 E E_\xi\|\widehat Z_1^\xi\|_{\infty,\mathcal I}^2
 \le Cn^{-1}.                                          \tag{19}
\]

It remains to prove the claimed covariance consistency.  First consider the un-interpolated processes.  Independence, (2), and the predictable isometry give the oracle covariance explicitly as
\[
 \Sigma_1(s,t)
 =\frac{s_1(s)s_1(t)}n
   \int_0^{s\wedge t}\frac{\lambda_1(u)}{H_D(u)}\,du.  \tag{20}
\]
Because each patient has at most one failure, the conditional covariance of the patient multiplier is
\[
 \widehat\Sigma_1(s,t)
 =\widehat S_D^{\mathrm{PL}}(s-)
  \widehat S_D^{\mathrm{PL}}(t-)
  \int_0^{s\wedge t}\frac{J_Y(u)}{Y(u)^2}\,dN(u),     \tag{21}
\]
up to the immaterial choice of left-continuous representatives at deterministic grid points.  Define
\[
 \mathcal A_n(v)=\int_0^v\frac{nJ_Y(u)}{Y(u)^2}\,dN(u),
 \qquad
 \mathcal A(v)=\int_0^v\frac{\lambda_1(u)}{H_D(u)}\,du. \tag{22}
\]
On \(A_n\), substitution of \(dN=dM+Y\lambda_1du\) gives
\[
 \mathcal A_n(v)-\mathcal A(v)
 =\int_0^v\frac n{Y(u)^2}\,dM(u)
  +\int_0^v\left\{\frac n{Y(u)}-\frac1{H_D(u)}\right\}
       \lambda_1(u)\,du.                               \tag{23}
\]
Stopping at \(\sigma_n\), the squared maximal expectation of the first term is at most
\[
 4E\int_0^{\sigma_n}\frac{n^2}{Y(u)^4}
       \,d\langle M\rangle(u)
 =4E\int_0^{\sigma_n}\frac{n^2\lambda_1(u)}{Y(u)^3}\,du
 \le Cn^{-1}.                                          \tag{24}
\]
For the second term, inversion is Lipschitz on
\([h_0/2,\infty)\), so (4) implies
\[
 E\left[\mathbf 1_{A_n}
 \sup_{v\in\mathcal I}
 \left|\int_0^v
 \left\{\frac n{Y(u)}-\frac1{H_D(u)}\right\}
 \lambda_1(u)\,du\right|^2\right]
 \le Cn^{-1}.                                          \tag{25}
\]
Equations (3) and (23)--(25) show, uniformly over the class,
\[
 \sup_{v\in\mathcal I}|\mathcal A_n(v)-\mathcal A(v)|=o_P(1). \tag{26}
\]
Equations (3), (6), and (15)--(17) also give
\[
 \|\widehat S_D^{\mathrm{PL}}-s_1\|_\infty=o_P(1)    \tag{27}
\]
uniformly.  Multiplying (20)--(21) by \(n\), then using (22), (26), (27), and the fixed bound on \(\mathcal A(\bar\tau)\), yields
\[
 n\|\widehat\Sigma_1-\Sigma_1\|_{\infty,\mathcal I^2}
 \longrightarrow0                                     \tag{28}
\]
in probability uniformly over \(\mathfrak P_1\).  Bilinear interpolation of covariance kernels is again a supremum-norm contraction, so (28) is exactly the assertion for the interpolated processes in the statement.

At time zero all counting-process integrals in (7), (18), (20), and (21) vanish.  The proof nowhere divides by a coordinate variance, so it also permits any other zero-variance coordinate.  No conditional quantile or coverage assertion has been used.  Every displayed bound used only the common hazard, the fixed overlap and smoothness constants, and independence; it is therefore uniform despite donor-specific censoring. \(\square\)